\subsection{Module 1 --- Geometry setup}
\label{sec:code-geometry}

This module builds, once per simulation, all coordinates needed by the induction system: the spanwise distribution of blade nodes, the bound-vortex segments at the $c/4$ line, the control points at the $3c/4$ line, and the helical wake nodes for every blade, every radial station, and every azimuthal step. The blade is divided into $N$ segments using either a uniform or a cosine spanwise distribution,
\begin{align}
r_i^{\text{unif}}   &= r_{\text{root}} + (R-r_{\text{root}})\,\frac{i}{N}, \label{eq:unif-disc}\\
r_i^{\text{cos}}    &= r_{\text{root}} + \tfrac{1}{2}(R-r_{\text{root}})\bigl(1-\cos(\pi i/N)\bigr), \label{eq:cos-disc}
\end{align}
for $i=0,\dots,N$. The cosine distribution clusters segments near root and tip and is preferred when tip-vortex resolution dominates the error budget. The wake nodes are placed using equations~\eqref{eq:wake-x}--\eqref{eq:U-wake} for a parameter sweep $t_k = k\,\Delta\psi/\Omega$, $k=0,\dots,N_{\text{az}}$, where $N_{\text{az}}=N_{\text{wake}}\cdot 2\pi/\Delta\psi$ is the total number of azimuthal wake stations.

\begin{lstlisting}[language=Python]
def build_geometry(
    R: float, r_root: float, N: int,
    N_blades: int, N_wake: int, dpsi_deg: float,
    U_inf: float, Omega: float, a_w: float,
    distribution: str = "cosine",
) -> dict:
    """Return a dict with bound nodes, control points,
    and wake node coordinates for all blades."""
    ...
\end{lstlisting}

The default values used throughout this report are listed in Table~\ref{tab:geom-params}.

\begin{table}[H]
\centering
\caption{Geometry and discretisation parameters.}
\label{tab:geom-params}
\begin{tabular}{@{}llcl@{}}
\toprule
Symbol & Parameter & Default & Role \\ \midrule
$N$            & spanwise segments       & 30          & resolution of bound circulation \\
$N_{\text{wake}}$ & wake rotations       & 10          & truncation length of helical wake \\
$\Delta\psi$   & azimuthal step          & $10^\circ$  & resolution of helical wake \\
$a_w$          & assumed wake induction  & 0.25        & convection speed in~\eqref{eq:U-wake} \\
$r_{\text{core}}$ & vortex core radius   & $10^{-3}R$  & regularises~\eqref{eq:katz-plotkin} \\
distribution   & spanwise law            & cosine      & uniform~\eqref{eq:unif-disc} or cosine~\eqref{eq:cos-disc} \\
\bottomrule
\end{tabular}
\end{table}

\subsection{Module 2 --- Biot--Savart kernel}
\label{sec:code-biot-savart}

The kernel is the lowest-level routine of the code: given a single straight filament defined by its endpoints $\vec{X}_1,\vec{X}_2$ and circulation $\Gamma$, it returns the velocity it induces at a control point $\vec{X}_p$. The implementation follows directly from the Katz \& Plotkin formula~\eqref{eq:katz-plotkin}, factored through the scalar $K$ in~\eqref{eq:K-scalar} so that the three components are obtained as in~\eqref{eq:UVW-components}.

\begin{lstlisting}[language=Python]
def filament_velocity(
    X1: np.ndarray, X2: np.ndarray, XP: np.ndarray,
    GAMMA: float, r_core: float,
) -> np.ndarray:
    """Velocity (u, v, w) at XP induced by a straight
    vortex filament X1->X2 of circulation GAMMA.
    XP can be (3,) or (M, 3) for vectorised evaluation."""
    R1 = XP - X1
    R2 = XP - X2
    cross = np.cross(R1, R2)
    cross_sq = np.einsum("...i,...i->...", cross, cross)
    # Core regularisation
    perp = np.sqrt(cross_sq) / np.linalg.norm(X2 - X1)
    inside_core = perp < r_core
    # K factor and (u, v, w) elsewhere
    K = (GAMMA / (4 * np.pi * cross_sq)) * (
        np.einsum("i,...i->...", X2 - X1, R1) / np.linalg.norm(R1, axis=-1)
      - np.einsum("i,...i->...", X2 - X1, R2) / np.linalg.norm(R2, axis=-1)
    )
    uvw = K[..., None] * cross
    uvw[inside_core] = 0.0
    return uvw
\end{lstlisting}

The core check triggers whenever the perpendicular distance from $\vec{X}_p$ to the filament drops below $r_{\text{core}}$; inside that radius the induced velocity is overwritten with zero, suppressing the $1/|\vec{R}_1\times\vec{R}_2|^2$ singularity of~\eqref{eq:K-scalar}. Because this function is called $\sim N^2 N_{\text{az}} N_b$ times during the assembly of Module~3, every operation is written as a vectorised NumPy expression that broadcasts over the leading axes of \texttt{XP}; no Python-level loops are used inside it.

\subsection{Module 3 --- Influence matrix assembly}
\label{sec:code-influence}

This module fills the three $N\times N$ matrices $[\mathbf{U}],[\mathbf{V}],[\mathbf{W}]$ entering~\eqref{eq:influence-matrix}. The algorithm is a double loop over control points $i=1,\dots,N$ and \emph{horseshoe rings} $j=1,\dots,N$. One horseshoe ring associated with the $j$-th spanwise segment of one reference blade contains: (i) the bound segment between the $j$-th and $(j{+}1)$-th $c/4$ nodes, and (ii) the two trailing legs leaving those nodes, each represented as a chain of $N_{\text{az}}$ straight segments following the helix~\eqref{eq:wake-x}--\eqref{eq:wake-z}. To enforce the rotor periodicity, the same ring is replicated and rotated to all $N_b$ blades, and the contributions are summed.

\begin{lstlisting}[language=Python]
def assemble_influence(
    cp: np.ndarray,        # (N, 3) control points on reference blade
    bound: np.ndarray,     # (N+1, 3) bound nodes on reference blade
    wake: np.ndarray,      # (N_b, N+1, N_az+1, 3) helical wake nodes
    r_core: float,
) -> tuple[np.ndarray, np.ndarray, np.ndarray]:
    N = cp.shape[0]
    U = np.zeros((N, N)); V = np.zeros((N, N)); W = np.zeros((N, N))
    for j in range(N):
        # --- bound segment of horseshoe j (reference blade only)
        # --- trailing leg from node j  (all blades, all azimuths)
        # --- trailing leg from node j+1 (all blades, all azimuths,
        #     opposite circulation sign)
        # Each call to filament_velocity is unit-strength (GAMMA = 1)
        # and accumulated into U[:, j], V[:, j], W[:, j].
        ...
    return U, V, W
\end{lstlisting}

Because the wake is frozen, the matrices depend only on the prescribed geometry. They are computed once at the start of the simulation and reused at every iteration of the non-linear solver in Module~5: the per-iteration cost of obtaining the induced velocities is reduced from a full Biot--Savart sweep ($\mathcal{O}(N^2 N_{\text{az}} N_b)$ kernel calls) to a single matrix--vector multiply~\eqref{eq:influence-matrix} of cost $\mathcal{O}(N^2)$, which is the main reason the frozen-wake assumption is computationally attractive.

\subsection{Module 4 --- Airfoil polar interpolation}
\label{sec:code-polar}

The 2D airfoil data for the DU 95-W-180 are loaded from the supplied \texttt{polar DU95W180.xlsx} into NumPy arrays of $\alpha$, $C_l$, $C_d$ tabulated values. The lift and drag coefficients entering~\eqref{eq:lift}--\eqref{eq:drag} are obtained by linear interpolation,
\begin{lstlisting}[language=Python]
def airfoil_coefficients(
    alpha: np.ndarray, polar: dict
) -> tuple[np.ndarray, np.ndarray]:
    Cl = np.interp(alpha, polar["alpha"], polar["Cl"])
    Cd = np.interp(alpha, polar["alpha"], polar["Cd"])
    return Cl, Cd
\end{lstlisting}
The \texttt{numpy.interp} routine extrapolates with the boundary values of the table (flat extrapolation) outside the tabulated range, which for the DU 95-W-180 covers approximately $\alpha\in[-15^\circ,30^\circ]$. Operating points venturing outside this window are therefore approximated rather than physically modelled, and a flag is raised in the post-processing if any control point exceeds the polar range. This module is reused unchanged from the BEM assignment.

\subsection{Module 5 --- Iteration loop}
\label{sec:code-iteration}

Module~5 implements Algorithm~\ref{alg:lifting-line}. With the influence matrices already in memory, each iteration consists of a matrix--vector multiply~\eqref{eq:influence-matrix}, the algebra of equations~\eqref{eq:V-axial}--\eqref{eq:phi-alpha}, a polar lookup, and a Kutta--Joukowski update~\eqref{eq:gamma-closure} with under-relaxation.

\begin{lstlisting}[language=Python]
def solve_circulation(
    U_mat, V_mat, W_mat, geom, polar,
    U_inf, Omega, w_relax=0.3, tol=1e-4, max_iter=1000,
) -> dict:
    Gamma = np.zeros(geom["N"])
    for k in range(max_iter):
        # (1) induced velocities at all control points
        u_ind = U_mat @ Gamma
        v_ind = V_mat @ Gamma
        w_ind = W_mat @ Gamma

        # (2) velocity triangle  -> eqs (V_axial), (V_tan), (V_p)
        V_ax  = U_inf + u_ind                       # axial component
        V_tan = Omega * geom["r"] + tangential(v_ind, w_ind, geom)
        V_p   = np.hypot(V_ax, V_tan)

        # (3) angle of attack    -> eq (phi-alpha)
        phi   = np.arctan2(V_ax, V_tan)
        alpha = phi - geom["twist"] - geom["pitch"]

        # (4) polar lookup       -> Module 4
        Cl, Cd = airfoil_coefficients(np.degrees(alpha), polar)

        # (5) new circulation    -> eq (gamma-closure)
        Gamma_new = 0.5 * geom["chord"] * V_p * Cl

        # (6) under-relaxation
        Gamma_next = (1 - w_relax) * Gamma + w_relax * Gamma_new

        # (7) convergence check
        denom = max(np.max(np.abs(Gamma_next)), 1e-12)
        if np.max(np.abs(Gamma_next - Gamma)) / denom < tol:
            Gamma = Gamma_next
            break
        Gamma = Gamma_next
    return {"Gamma": Gamma, "alpha": alpha, "phi": phi,
            "u_ind": u_ind, "v_ind": v_ind, "w_ind": w_ind, "iters": k}
\end{lstlisting}

For the baseline configuration the loop converges to $\epsilon=10^{-4}$ in roughly $80$--$150$ iterations. Two failure modes have been observed: (i) for $w\gtrsim 0.5$ the residual oscillates rather than decays, especially near $\lambda=10$ where the system is stiffest; (ii) for $w\lesssim 0.1$ convergence is monotonic but the iteration count grows beyond a thousand. The default $w=0.3$ is a robust compromise across the three tip-speed ratios required by the assignment.

\subsection{Module 6 --- Post-processing and sensitivity study}
\label{sec:code-postprocess}

After convergence, the spanwise distributions are integrated to obtain the integral rotor performance. The sectional axial and azimuthal forces from~\eqref{eq:F-axial}--\eqref{eq:F-azim} are integrated using the trapezoid rule, multiplied by the number of blades, and converted to thrust and torque,
\begin{align}
T &= N_b\int_{r_{\text{root}}}^{R} F_{\text{axial}}(r)\,\mathrm{d}r,
& Q &= N_b\int_{r_{\text{root}}}^{R} r\,F_{\text{azim}}(r)\,\mathrm{d}r,
& P = \Omega\,Q.
\label{eq:T-Q-integration}
\end{align}
Equations~\eqref{eq:CT}--\eqref{eq:CP} then yield $C_T$ and $C_P$, while the induction factors are reconstructed from the converged induced velocities through~\eqref{eq:induction-factors}. The sensitivity study is implemented as a wrapper that re-runs the full pipeline of Modules~1--5 for each value of one parameter while keeping the others at their defaults; the parameters and ranges are listed in Table~\ref{tab:sensitivity}.

\begin{lstlisting}[language=Python]
def sensitivity_sweep(
    parameter: str, values: list, defaults: dict
) -> list[dict]:
    results = []
    for val in values:
        cfg = {**defaults, parameter: val}
        geom = build_geometry(**cfg)
        U_m, V_m, W_m = assemble_influence(**geom_args(geom))
        out = solve_circulation(U_m, V_m, W_m, geom, polar,
                                cfg["U_inf"], cfg["Omega"])
        results.append({"value": val, **post_process(out, geom, cfg)})
    return results
\end{lstlisting}

\begin{table}[H]
\centering
\caption{Sensitivity-study parameters and assignment mapping.}
\label{tab:sensitivity}
\begin{tabular}{@{}llccl@{}}
\toprule
Parameter & Symbol & Default & Sweep range & Question addressed \\ \midrule
wake convection      & $a_w$            & 0.25        & $[0,\,0.5]$                         & convection speed sensitivity \\
spanwise resolution  & $N$              & 30          & $\{5,10,20,40,80\}$                 & blade discretisation (also uniform vs cosine) \\
wake length          & $N_{\text{wake}}$ & 10         & $\{1,3,5,10,20\}$                   & convergence with wake length \\
azimuthal step       & $\Delta\psi$     & $10^\circ$  & $\{5^\circ,10^\circ,20^\circ\}$     & azimuthal discretisation \\
\bottomrule
\end{tabular}
\end{table}

The full set of output quantities produced by the post-processing pipeline, the equation used to compute each, and its non-dimensionalisation, are summarised in Table~\ref{tab:outputs}.

\begin{table}[H]
\centering
\caption{Output quantities produced by the lifting line solver.}
\label{tab:outputs}
\begin{tabular}{@{}llcl@{}}
\toprule
Quantity & Symbol & Computed via & Non-dimensional form \\ \midrule
inflow angle           & $\Phi$         & \eqref{eq:phi-alpha}                & dimensionless (rad / $^\circ$) \\
angle of attack        & $\alpha$       & \eqref{eq:phi-alpha}                & dimensionless (rad / $^\circ$) \\
axial induction        & $a$            & \eqref{eq:induction-factors}        & dimensionless \\
tangential induction   & $a'$           & \eqref{eq:induction-factors}        & dimensionless \\
bound circulation      & $\Gamma$       & \eqref{eq:gamma-closure}            & $\hat{\Gamma}$~\eqref{eq:gamma-hat} \\
axial force per span   & $F_{\text{axial}}$ & \eqref{eq:F-axial}              & $\hat{F}$~\eqref{eq:F-hat} \\
azimuthal force per span & $F_{\text{azim}}$ & \eqref{eq:F-azim}             & $\hat{F}$~\eqref{eq:F-hat} \\
thrust                 & $T$            & \eqref{eq:T-Q-integration}          & $C_T$~\eqref{eq:CT} \\
power                  & $P=\Omega Q$   & \eqref{eq:T-Q-integration}          & $C_P$~\eqref{eq:CP} \\
\bottomrule
\end{tabular}
\end{table}